\documentclass[11pt]{article}
\usepackage[letterpaper,margin=0.83in,headheight=16pt]{geometry}
\usepackage[T1]{fontenc}
\usepackage[utf8]{inputenc}
\usepackage{lmodern}
\usepackage{microtype}
\usepackage{amsmath,amssymb}
\usepackage{xcolor}
\usepackage{hyperref}
\usepackage{enumitem}
\usepackage{booktabs,longtable,tabularx,array}
\usepackage{fancyhdr}
\usepackage{titlesec}
\usepackage{listings}
\usepackage{tikz}
\usepackage{float}
\definecolor{suiteblue}{HTML}{17365D}
\definecolor{suitelight}{HTML}{EAF1F8}
\definecolor{suitegray}{HTML}{556070}
\hypersetup{colorlinks=true,linkcolor=suiteblue,urlcolor=suiteblue,breaklinks=true}
\titleformat{\section}{\Large\bfseries\color{suiteblue}}{\thesection}{0.55em}{}
\titleformat{\subsection}{\large\bfseries\color{suiteblue}}{\thesubsection}{0.55em}{}
\setlist[itemize]{itemsep=0.25em,topsep=0.35em,leftmargin=1.4em}
\setlist[enumerate]{itemsep=0.25em,topsep=0.35em,leftmargin=1.6em}
\setlength{\parindent}{0pt}
\setlength{\parskip}{0.47em}
\emergencystretch=2em
\lstset{basicstyle=\ttfamily\footnotesize,breaklines=true,frame=single,columns=fullflexible,showstringspaces=false}
\pagestyle{fancy}\fancyhf{}
\lhead{Game Design \& Engineering}\rhead{Workbook Template}\cfoot{\thepage}
\begin{document}
{\huge\bfseries\color{suiteblue} Mechanics Modeling Workbook -- Template}\par
\vspace{0.15in}
{\normalsize\color{suitegray} Fictional example annotations plus reusable analytical prompts}\par
\vspace{0.15in}\hrule\vspace{0.15in}

\section{Record control}
Project: [Project]. Revision: [Revision]. Practitioner: [Name]. Evidence repository: [URI].


\section{Evidence discipline}
For each exercise, distinguish an assumption, result established by computation, result established by player observation, and an interpretation. The example entries are instructional and are not observations from a real game.


\section{Formalize the state}
Exercise prompt: Define S with current health in 0..10, energy in 0..5, and a turn flag. Enumerate at least three legal actions and five illegal actions.

Expected deliverable: State schema; transition preconditions; four invariants; transition table.


\subsection{Hypothesis}
Illustrative example: A small design change affects a specified, observable player action.

Project-specific response: [Document the evidence and rationale.]


\subsection{Assumptions}
Illustrative example: The simulated population and state constraints are explicitly defined.

Project-specific response: [Document the evidence and rationale.]


\subsection{Model or procedure}
Illustrative example: Use a repeatable implementation and record input parameters.

Project-specific response: [Document the evidence and rationale.]


\subsection{Results and limitations}
Illustrative example: Quantitative result is supported by a run log; enjoyment remains unverified.

Project-specific response: [Document the evidence and rationale.]


\subsection{Decision and next test}
Illustrative example: Keep, revise, or retire the feature after a new independent study.

Project-specific response: [Document the evidence and rationale.]


\section{Check conservation}
Exercise prompt: Specify a spend action, a recharge action, and a failure to spend. Show that spent energy is never negative.

Expected deliverable: Invariant proof sketch and ten executable test cases.


\subsection{Hypothesis}
Illustrative example: A small design change affects a specified, observable player action.

Project-specific response: [Document the evidence and rationale.]


\subsection{Assumptions}
Illustrative example: The simulated population and state constraints are explicitly defined.

Project-specific response: [Document the evidence and rationale.]


\subsection{Model or procedure}
Illustrative example: Use a repeatable implementation and record input parameters.

Project-specific response: [Document the evidence and rationale.]


\subsection{Results and limitations}
Illustrative example: Quantitative result is supported by a run log; enjoyment remains unverified.

Project-specific response: [Document the evidence and rationale.]


\subsection{Decision and next test}
Illustrative example: Keep, revise, or retire the feature after a new independent study.

Project-specific response: [Document the evidence and rationale.]


\section{Explore state-space explosion}
Exercise prompt: Compute the naive upper bound on states. Reduce it by invariant-based constraints; explain why the remaining bound is not the reachable-state count.

Expected deliverable: Counting argument; bounded exhaustive test plan.


\subsection{Hypothesis}
Illustrative example: A small design change affects a specified, observable player action.

Project-specific response: [Document the evidence and rationale.]


\subsection{Assumptions}
Illustrative example: The simulated population and state constraints are explicitly defined.

Project-specific response: [Document the evidence and rationale.]


\subsection{Model or procedure}
Illustrative example: Use a repeatable implementation and record input parameters.

Project-specific response: [Document the evidence and rationale.]


\subsection{Results and limitations}
Illustrative example: Quantitative result is supported by a run log; enjoyment remains unverified.

Project-specific response: [Document the evidence and rationale.]


\subsection{Decision and next test}
Illustrative example: Keep, revise, or retire the feature after a new independent study.

Project-specific response: [Document the evidence and rationale.]


\section{Analyze emergent actions}
Exercise prompt: Combine pushing, collision, and resource use to construct two unplanned but valid strategies.

Expected deliverable: Counterexample traces and design review.


\subsection{Hypothesis}
Illustrative example: A small design change affects a specified, observable player action.

Project-specific response: [Document the evidence and rationale.]


\subsection{Assumptions}
Illustrative example: The simulated population and state constraints are explicitly defined.

Project-specific response: [Document the evidence and rationale.]


\subsection{Model or procedure}
Illustrative example: Use a repeatable implementation and record input parameters.

Project-specific response: [Document the evidence and rationale.]


\subsection{Results and limitations}
Illustrative example: Quantitative result is supported by a run log; enjoyment remains unverified.

Project-specific response: [Document the evidence and rationale.]


\subsection{Decision and next test}
Illustrative example: Keep, revise, or retire the feature after a new independent study.

Project-specific response: [Document the evidence and rationale.]


\section{Review decision}
Reviewer: [Name]. Findings: [Finding IDs]. Status: [Draft, revision needed, or accepted].


\end{document}
